\documentclass[11pt,a4paper]{article}

\usepackage[utf8]{inputenc}
\usepackage[T1]{fontenc}
\usepackage[polish]{babel}
\usepackage{geometry}
\usepackage{booktabs}
\usepackage{tabularray}
\usepackage{afterpage}
\usepackage{rotating}
\usepackage{longtable}
\usepackage{array}
\usepackage{xcolor}
\usepackage{hyperref}
\usepackage{titlesec}
\usepackage{parskip}
\usepackage{enumitem}
\usepackage{fancyhdr}
\usepackage{pdflscape}

\geometry{margin=2.5cm}

\definecolor{accent}{RGB}{30, 60, 100}

\titleformat{\section}{\large\bfseries\color{accent}}{}{0em}{}[\titlerule]
\titleformat{\subsection}{\normalsize\bfseries}{}{0em}{}
\titleformat{\subsubsection}{\normalsize\itshape}{}{0em}{}

\pagestyle{fancy}
\fancyhf{}
\rhead{\small Specyfikacja Projektu}
\lhead{\small Projekt: \textit{Feedback}}
\cfoot{\thepage}
\renewcommand{\headrulewidth}{0.4pt}

\hypersetup{colorlinks=true, linkcolor=accent, urlcolor=accent}

\begin{document}

% ── STRONA TYTUŁOWA ──────────────────────────────────────────────
\begin{titlepage}
\centering
\vspace*{2cm}
{\LARGE\bfseries\color{accent} Specyfikacja Projektu}\\[0.5cm]
{\large Projekt: \textit{Feedback}}\\[0.3cm]
{\normalsize Mechanika parowania jako klucz do sukcesu}\\[2cm]

\begin{tabular}{ll}
\textbf{Kierownik projektu:}   & Kacper Trzciński \\[0.3em]
\textbf{Przewodnicząca KS:}    & Aleksandra Sobótka \\[0.3em]
\textbf{Główny Dostawca:}      & Aleksander Płocharski \\[0.3em]
\textbf{Główny Użytkownik:}    & Tomasz Herman \\[0.3em]
\textbf{Zespół:}               & A. Bielicki, R. Głasek, J. Pączka, \\
                               & M. Popkowicz, F. Shramko, K. Trzciński \\[0.3em]
\textbf{Data opracowania:}     & \today \\[0.3em]
\end{tabular}

\vfill
{\small Zarządzanie Przedsięwzięciami Informatycznymi\\
Wydział Matematyki i Nauk Informacyjnych, Politechnika Warszawska}
\end{titlepage}

\newpage
\tableofcontents
\newpage

\section{Wprowadzenie}

\subsection{Krótki opis gry}
„Feedback” to innowacyjna gra typu \textit{puzzle-action platformer}, która redefiniuje pojęcie walki w grach 2D. W przeciwieństwie do klasycznych przedstawicieli gatunku, gdzie dominują systemy oparte na zadawaniu obrażeń (agresji), „Feedback” eliminuje tradycyjną broń ofensywną na rzecz \textbf{mechaniki parowania (parry)}.

Główne założenia rozgrywki opierają się na trzech filarach:
\begin{itemize}[noitemsep]
    \item \textbf{Defensywa jako jedyna forma interakcji:} Gracz nie atakuje wprost. Każde pokonanie przeciwnika wymaga przejęcia jego siły – poprzez odbicie pocisku lub zneutralizowanie ataku fizycznego w idealnym momencie.
    \item \textbf{Manipulacja fizyką i momentum:} System parowania jest zintegrowany z modelem poruszania się. Udane parowanie pozwala graczowi na zmianę własnego wektora ruchu, co umożliwia pokonywanie przeszkód terenowych (np. \textit{rocket jump} po sparowaniu wybuchu).
    \item \textbf{Zagadki środowiskowe:} Poziomy są zaprojektowane jako logiczne labirynty, w których przeciwnicy pełnią rolę „generatorów energii”. Aby otworzyć przejście, gracz musi sprowokować atak i precyzyjnie przekierować go w odpowiedni element otoczenia.
\end{itemize}

\subsection{Cel projektu}
Nadrzędnym celem projektu jest stworzenie grywalnego prototypu, który udowodni, że mechanika obronna może być fundamentem dynamicznej i satysfakcjonującej pętli rozgrywki. Cele operacyjne obejmują:
\begin{itemize}[noitemsep]
    \item \textbf{Implementacja logiki systemowej parowania:} Opracowanie skryptów obsługujących tzw. \textit{frame-data}, czyli precyzyjne okna czasowe (\textit{parry windows}), w których akcja gracza kończy się sukcesem.
    \item \textbf{Opracowanie systemu transferu energii:} Wykorzystanie wbudowanych mechanizmów silnika fizycznego do przeliczania i zmiany momentum postaci po udanym bloku.
    \item \textbf{Przygotowanie trzech poziomów testowych:} Stworzenie etapów o narastającym stopniu skomplikowania zagadek logicznych.
    \item \textbf{Dostarczenie pełnej dokumentacji:} Przygotowanie artefaktów zarządczych zgodnych z metodyką PRINCE2.
\end{itemize}

\subsection{Kontekst i innowacyjność}
Projekt powstaje jako odpowiedź na stagnację w gatunku gier akcji, gdzie większość innowacji ogranicza się do modyfikacji statystyk postaci. „Feedback” przesuwa akcent na \textbf{czyste umiejętności techniczne gracza (skill-based design)}. Zamiast zbierania przedmiotów zwiększających siłę ataku, gracz rozwija własną biegłość w obsłudze jednej, doskonale dopracowanej mechaniki. Innowacja polega na transformacji parowania z narzędzia czysto pomocniczego (znanego z serii \textit{Souls} czy \textit{Sekiro}) w główny język interakcji ze światem gry.

\section{Użytkownicy}

\subsection{Grupa docelowa}
Głównym odbiorcą gry są gracze preferujący gry zręcznościowe połączone z wyzwaniami logicznymi, którzy szukają innowacyjnych mechanik rozgrywki wykraczających poza standardowe schematy walki.

\subsection{Charakterystyka użytkownika}
Idealny użytkownik „Feedback” to osoba, która:
\begin{itemize}[noitemsep]
    \item Ceni wysoką responsywność sterowania i dynamikę ruchu.
    \item Szuka gier wymagających analizy zachowań przeciwników i strategicznego planowania działań obronnych.
    \item Potrafi czerpać satysfakcję z opanowania technicznych mechanik gry (tzw. \textit{skill-based gameplay}).
\end{itemize}

\subsection{Odbiorca wtórny (Interesariusze)}
Zgodnie ze strukturą PRINCE2, istotnymi odbiorcami efektów projektu są:
\begin{itemize}[noitemsep]
    \item \textbf{Główny Użytkownik:} Tomasz Herman, reprezentujący perspektywę końcowego odbiorcy i weryfikujący grywalność mechanik.
    \item \textbf{Komitet Sterujący:} Nadzorujący zgodność projektu z celami dydaktycznymi przedmiotu ZPI.
    \item \textbf{Przyszli studenci:} Jako odbiorcy końcowej prezentacji i demo gry zaplanowanej na semestr letni 2027.
\end{itemize}

\section{Opis gry}

\textit{Feedback} jest dwuwymiarową grą platformowo--zręcznościową z elementami logicznymi, której głównym założeniem projektowym jest całkowite odejście od tradycyjnych systemów walki opartych na zadawaniu obrażeń. Fundament rozgrywki stanowi mechanika parowania (\textit{parry}), pełniąca jednocześnie funkcję defensywną, ruchową oraz środowiskową.

Gracz wciela się w postać poruszającą się po serii zamkniętych poziomów o strukturze logicznych torów przeszkód. Każdy poziom wymaga nie tylko zręczności i precyzyjnego sterowania, ale również analizy otoczenia oraz zrozumienia zależności pomiędzy rozmieszczeniem przeciwników, przeszkód i elementów interaktywnych.

Rozgrywka opiera się na minimalistycznym zestawie mechanik:
\begin{itemize}[noitemsep]
    \item poruszaniu się,
    \item skakaniu,
    \item wykonywaniu parowania.
\end{itemize}

Pomimo ograniczonej liczby dostępnych akcji, gra oferuje dużą różnorodność sytuacji dzięki projektowi poziomów oraz interakcjom opartym na fizyce. Udane parowanie pozwala nie tylko neutralizować zagrożenia, ale również przejmować energię przeciwników i wykorzystywać ją do manipulowania własnym ruchem lub wpływania na otoczenie. Mechanika ta umożliwia m.in. wykonywanie dalekich skoków, zmianę kierunku ruchu w powietrzu czy aktywowanie elementów środowiska za pomocą odbitych pocisków.

Przeciwnicy pełnią w grze podwójną rolę — stanowią zagrożenie dla gracza, ale jednocześnie są integralną częścią zagadek środowiskowych. W wielu przypadkach ukończenie poziomu wymaga świadomego sprowokowania ataku przeciwnika, a następnie odpowiedniego przekierowania jego efektu w celu odblokowania dalszej drogi.

Projekt poziomów został podporządkowany idei stopniowego rozwijania kompetencji gracza. Kolejne etapy wprowadzają nowe elementy interakcji i coraz bardziej wymagające kombinacje mechanik, zachęcając do eksperymentowania oraz doskonalenia umiejętności. Gra premiuje zarówno refleks i precyzję wykonania, jak i zdolność planowania optymalnej ścieżki przejścia.

Istotnym elementem doświadczenia jest szybki cykl prób i błędów. Niepowodzenie skutkuje natychmiastowym restartem poziomu, dzięki czemu gracz może płynnie podejmować kolejne próby bez zbędnych przerw. Rozwiązanie to wspiera dynamiczny charakter rozgrywki i wzmacnia satysfakcję wynikającą z opanowywania trudnych sekwencji ruchu.

\section{Pętla rozgrywki}

Poniższa sekcja opisuje poszczególne etapy pętli rozgrywki gry.
 \subsection{Gameplay loop}

Pojedynczy poziom gry \textit{Feedback} składa się z trzech podstawowych kategorii elementów:
\begin{itemize}
    \item \textbf{Obiekty statyczne} --- ściany, podłogi oraz nieruchome platformy tworzące podstawową strukturę poziomu.
    \item \textbf{Obiekty dynamiczne} --- ruchome platformy, trampoliny, klucze, zamknięte drzwi oraz inne interaktywne elementy środowiska posiadające określone właściwości fizyczne lub logiczne.
    \item \textbf{Przeciwnicy} --- główne źródło zagrożenia, ale jednocześnie integralny element zagadek środowiskowych i mechanik poziomu.
\end{itemize}

Celem gracza jest dotarcie od punktu początkowego poziomu do wyznaczonego miejsca końcowego, np. widocznego wyjścia lub aktywowanego portalu. Osiągnięcie tego punktu oznacza ukończenie poziomu i odblokowanie kolejnego etapu rozgrywki.

Każdy element otoczenia może wpływać na przebieg rozgrywki. Nawet pasywne obiekty, takie jak ściany czy podłogi, pełnią istotną funkcję w systemie ruchu i fizyki. Odpowiednio ustawione powierzchnie umożliwiają wykonywanie odbić pod określonym kątem, co pozwala graczowi zachować lub zmienić momentum ruchu.

Kluczowym aspektem gry jest struktura poziomów. Gracz dysponuje ograniczonym zestawem akcji --- poruszaniem się, skokiem oraz parowaniem --- jednak różnorodność wyzwań wynika z konfiguracji przeszkód, przeciwników i interaktywnych elementów środowiska. Poszczególne poziomy są projektowane w taki sposób, aby wymuszać kreatywne wykorzystanie tych samych mechanik w nowych kontekstach.

Ze względu na logiczno--zręcznościowy charakter gry, odnalezienie poprawnej ścieżki przejścia może wymagać wielu prób oraz analizy zachowania przeciwników i otoczenia. Proces rozwiązywania poziomu opiera się na metodzie prób i błędów, obserwacji zależności pomiędzy elementami mapy oraz stopniowym doskonaleniu wykonania.

Po odkryciu optymalnego rozwiązania poziom powinien być możliwy do ukończenia w sposób płynny i dynamiczny. Gra kładzie nacisk na precyzję sterowania, refleks oraz satysfakcję wynikającą z perfekcyjnego wykonania wcześniej zaplanowanej sekwencji działań.

Próba ukończenia poziomu kończy się niepowodzeniem w momencie kontaktu postaci z zagrożeniem, np. przeciwnikiem, pociskiem lub przeszkodą środowiskową taką jak kolce. Po śmierci gracz natychmiast powraca do początku poziomu, co pozwala zachować wysoką dynamikę rozgrywki i zachęca do szybkiego podejmowania kolejnych prób.

\subsection{Progression loop}

Postęp w grze opiera się przede wszystkim na przechodzeniu kolejnych poziomów. Każdy ukończony etap odblokowuje następny, wprowadzający nowe konfiguracje przeszkód, przeciwników oraz mechanik środowiskowych.

Rozwój gracza nie jest związany z klasycznym systemem statystyk, ekwipunku ani odblokowywaniem nowych umiejętności postaci. Zestaw dostępnych akcji pozostaje niezmienny przez całą grę. Progresja opiera się zamiast tego na stopniowym rozwijaniu umiejętności samego gracza, takich jak refleks, precyzja ruchu, znajomość zachowań przeciwników oraz zdolność analizowania poziomów.

Kolejne etapy pełnią funkcję edukacyjną, prezentując nowe zastosowania istniejących mechanik. Początkowo elementy rozgrywki wprowadzane są w izolacji, a następnie stopniowo łączone w bardziej złożone układy wymagające szybkiego podejmowania decyzji oraz precyzyjnego wykonania.

Struktura progresji została zaprojektowana zgodnie z zasadą kumulatywnego zwiększania złożoności. Oznacza to, że wcześniej poznane mechaniki pozostają istotne również na późniejszych etapach gry, gdzie są wykorzystywane w nowych kontekstach i bardziej wymagających kombinacjach.

Pomimo ograniczonego nacisku na narrację, ukończenie ostatniego poziomu stanowi formalne zakończenie rozgrywki i główny cel gracza.

\subsection{Meta loop}

Ze względu na jednoosobowy i zręcznościowy charakter gry, system meta--progresji został celowo ograniczony. Projekt koncentruje się przede wszystkim na satysfakcji wynikającej z opanowywania mechanik oraz doskonalenia umiejętności gracza.

Jednocześnie struktura rozgrywki naturalnie sprzyja podejściu opartemu na poprawianiu wyników i optymalizacji przejść poziomów. Krótki czas trwania pojedynczych prób oraz wysoka precyzja mechanik zachęcają do wielokrotnego powracania do ukończonych etapów w celu uzyskania lepszych rezultatów.

Potencjalnym rozszerzeniem gry może być implementacja systemu rankingów oraz pomiaru czasu przejścia poziomów (\textit{time trial}). Rozwiązanie to umożliwiłoby porównywanie wyników pomiędzy graczami oraz wspierało społecznościowy aspekt rywalizacji bez naruszania podstawowych założeń projektowych.

Meta--progresja w \textit{Feedback} nie opiera się więc na wzmacnianiu postaci, lecz na rozwijaniu wiedzy o poziomach, doskonaleniu techniki oraz zwiększaniu efektywności wykonywania poszczególnych sekwencji działań.

\section{Wymagania funkcjonalne}
Wymagania funkcjonalne opisane zostały w formie historyjek użytkownika (\textit{user stories}) uzupełnionych kryteriami akceptacji wyrażonymi w notacji Gherkin (Mając--Kiedy--Wtedy). Takie połączenie pozwala zachować perspektywę odbiorcy oraz jednoznacznie zweryfikować spełnienie wymagania na etapie testów. Każdemu wymaganiu przypisano unikalny identyfikator (WF-XX) wykorzystywany w macierzy śledzenia oraz w rejestrze testów.

\subsection{Sterowanie postacią}

\subsubsection*{WF-01: Poruszanie się postaci}
\textbf{User story:} Jako gracz chcę poruszać postacią w lewo i w prawo, aby przemieszczać się po poziomie.

\textit{Kryteria akceptacji:}
\begin{itemize}[noitemsep]
    \item Mając aktywną postać na poziomie, kiedy gracz wciska klawisz ruchu w lewo lub w prawo, wtedy postać przyspiesza w danym kierunku do stałej prędkości bazowej.
    \item Mając postać poruszającą się w jednym kierunku, kiedy gracz puszcza klawisz ruchu, wtedy postać zatrzymuje się w ciągu jednej klatki (brak inercji w poziomie poza efektem parowania).
    \item Mając postać poruszającą się w jednym kierunku, kiedy gracz wciska klawisz ruchu w kierunku przeciwnym, wtedy zmiana kierunku następuje natychmiastowo.
\end{itemize}

\subsubsection*{WF-02: Skok}
\textbf{User story:} Jako gracz chcę móc skakać, aby pokonywać przeszkody pionowe i platformy.

\textit{Kryteria akceptacji:}
\begin{itemize}[noitemsep]
    \item Mając postać znajdującą się na powierzchni stałej, kiedy gracz wciska klawisz skoku, wtedy postać wykonuje skok o ustalonej wysokości bazowej.
    \item Mając postać w powietrzu w wyniku skoku, kiedy gracz ponownie wciska klawisz skoku, wtedy postać nie wykonuje kolejnego skoku.
    \item Mając postać w powietrzu, kiedy gracz wciska klawisz ruchu w lewo lub w prawo, wtedy postać reaguje na kierunek z ograniczonym wpływem na wektor prędkości (kontrola powietrzna).
\end{itemize}

\subsection{Mechanika parowania}

\subsubsection*{WF-03: Aktywacja parowania}
\textbf{User story:} Jako gracz chcę wykonywać parowanie, aby neutralizować ataki przeciwników i wykorzystywać ich energię.

\textit{Kryteria akceptacji:}
\begin{itemize}[noitemsep]
    \item Mając postać w dowolnym stanie (na stałym podłożu lub w powietrzu), kiedy gracz wciska klawisz parowania, wtedy postać przechodzi w stan parowania na zdefiniowane okno czasowe (\textit{parry window}).
    \item Mając postać w stanie parowania, kiedy do hitboxu parowania wchodzi atak (pocisk lub atak fizyczny), wtedy atak jest neutralizowany.
    \item Mając zakończone okno parowania, kiedy gracz ponownie wciska klawisz parowania, wtedy obowiązuje okres odnowienia (\textit{cooldown}) uniemożliwiający natychmiastowe powtórzenie akcji.
\end{itemize}

\subsubsection*{WF-04: Odbijanie pocisków}
\textbf{User story:} Jako gracz chcę sparowanym ciosem odbijać pociski przeciwników, aby przekierowywać je w przeciwników lub elementy otoczenia.

\textit{Kryteria akceptacji:}
\begin{itemize}[noitemsep]
    \item Mając postać w stanie parowania, kiedy pocisk przeciwnika trafia w hitbox parowania, wtedy pocisk zmienia kierunek na przeciwny do pierwotnego.
    \item Mając odbity pocisk, kiedy trafia on w przeciwnika, wtedy przeciwnik zostaje wyeliminowany.
    \item Mając odbity pocisk, kiedy trafia on w element interaktywny otoczenia (np. przełącznik), wtedy element zostaje aktywowany zgodnie ze swoją funkcją.
    \item Mając odbity pocisk, kiedy trafia on w postać gracza, wtedy postać ginie i poziom zostaje zrestartowany.
\end{itemize}

\subsubsection*{WF-05: Zmiana momentum po parowaniu}
\textbf{User story:} Jako gracz chcę zyskiwać dodatkowy ruch po udanym sparowaniu ataku fizycznego, aby pokonywać przepaście i osiągać niedostępne miejsca.

\textit{Kryteria akceptacji:}
\begin{itemize}[noitemsep]
    \item Mając postać w stanie parowania, kiedy atak fizyczny przeciwnika trafia w hitbox parowania, wtedy wektor pędu postaci zmienia się zgodnie z przekierowaniem energii ataku.
    \item Mając zmieniony wektor pędu, kiedy postać znajduje się w powietrzu, wtedy zachowuje ten pęd do momentu zetknięcia z powierzchnią lub kolejnego parowania.
    \item Mając postać po udanym parowaniu wybuchu, kiedy parowanie nastąpiło z dolnym hitboxem, wtedy postać uzyskuje pęd skierowany pionowo do góry (efekt \textit{rocket jump}).
\end{itemize}

\subsubsection*{WF-06: Wskazanie okna parowania}
\textbf{User story:} Jako gracz chcę otrzymywać natychmiastową informację zwrotną o wykonaniu parowania, aby uczyć się timingu mechaniki.

\textit{Kryteria akceptacji:}
\begin{itemize}[noitemsep]
    \item Mając wciśnięty klawisz parowania, kiedy okno parowania jest aktywne, wtedy postać wyświetla animację oraz efekt wizualny sygnalizujący gotowość obronną.
    \item Mając udane parowanie, kiedy atak zostaje zneutralizowany, wtedy gra wyświetla wyraźny efekt wizualny i akustyczny potwierdzenia.
    \item Mając nieudane parowanie (atak trafia poza oknem), kiedy postać zostaje trafiona, wtedy gra wyświetla efekt śmierci postaci.
\end{itemize}

\subsection{Przeciwnicy}

\subsubsection*{WF-07: Przeciwnik dystansowy}
\textbf{User story:} Jako gracz chcę zmierzyć się z przeciwnikiem strzelającym pociskami, aby wykorzystywać jego ataki jako element zagadki.

\textit{Kryteria akceptacji:}
\begin{itemize}[noitemsep]
    \item Mając przeciwnika dystansowego na poziomie, kiedy postać gracza znajduje się w jego polu widzenia, wtedy przeciwnik wystrzeliwuje pocisk w kierunku gracza w stałym interwale czasowym.
    \item Mając przeciwnika dystansowego, kiedy zostaje trafiony odbitym pociskiem, wtedy zostaje wyeliminowany z poziomu.
    \item Mając przeciwnika dystansowego, kiedy gracz nie znajduje się w jego polu widzenia, wtedy przeciwnik nie wystrzeliwuje pocisków.
\end{itemize}

\subsubsection*{WF-08: Przeciwnik fizyczny}
\textbf{User story:} Jako gracz chcę zmierzyć się z przeciwnikiem atakującym wręcz, aby sparować jego cios i wykorzystać przejęty pęd do ruchu.

\textit{Kryteria akceptacji:}
\begin{itemize}[noitemsep]
    \item Mając przeciwnika fizycznego, kiedy gracz znajduje się w jego zasięgu, wtedy przeciwnik wykonuje atak posiadający fazę zapowiedzi (\textit{telegraph}) oraz fazę aktywną.
    \item Mając atak przeciwnika w fazie zapowiedzi, kiedy gracz aktywuje parowanie tak, by okno parowania pokryło fazę aktywną, wtedy atak zostaje sparowany i przekazuje pęd graczowi.
    \item Mając parowanie nieudane, kiedy atak przeciwnika trafia w postać, wtedy postać ginie i poziom zostaje zrestartowany.
\end{itemize}

\subsection{Poziomy i progresja}

\subsubsection*{WF-09: Rozpoczęcie poziomu}
\textbf{User story:} Jako gracz chcę rozpocząć poziom z ustalonego punktu startowego, aby mieć powtarzalne warunki rozgrywki przy każdej próbie.

\textit{Kryteria akceptacji:}
\begin{itemize}[noitemsep]
    \item Mając wybrany poziom, kiedy poziom zostaje uruchomiony, wtedy postać pojawia się w zdefiniowanym dla tego poziomu punkcie startowym.
    \item Mając uruchomiony poziom, kiedy poziom rozpoczyna się, wtedy wszystkie elementy interaktywne (przełączniki, drzwi, przeciwnicy) znajdują się w stanie początkowym.
\end{itemize}

\subsubsection*{WF-10: Ukończenie poziomu}
\textbf{User story:} Jako gracz chcę móc ukończyć poziom przez dotarcie do punktu końcowego, aby przejść do kolejnego etapu.

\textit{Kryteria akceptacji:}
\begin{itemize}[noitemsep]
    \item Mając aktywny poziom, kiedy postać dociera do oznaczonego punktu końcowego, wtedy poziom zostaje zakończony.
    \item Mając ukończony poziom, kiedy zakończenie nastąpiło, wtedy gra automatycznie ładuje kolejny poziom z sekwencji lub wyświetla ekran końcowy, jeżeli był to ostatni poziom.
\end{itemize}

\subsubsection*{WF-11: Restart po śmierci}
\textbf{User story:} Jako gracz chcę szybko wracać do początku poziomu po śmierci, aby zachować dynamikę rozgrywki i podejmować kolejne próby bez zwłoki.

\textit{Kryteria akceptacji:}
\begin{itemize}[noitemsep]
    \item Mając aktywny poziom, kiedy postać gracza zostaje wyeliminowana (kontakt z przeciwnikiem, pociskiem, kolcami lub innym zagrożeniem), wtedy poziom zostaje zrestartowany w ciągu maksymalnie 1 sekundy.
    \item Mając zrestartowany poziom, kiedy restart się zakończył, wtedy stan poziomu i postaci odpowiada stanowi z momentu jego rozpoczęcia.
    \item Mając aktywny poziom, kiedy gracz wciska klawisz manualnego restartu, wtedy obowiązuje to samo zachowanie co przy śmierci postaci.
\end{itemize}

\subsubsection*{WF-12: Co najmniej trzy poziomy z zagadkami}
\textbf{User story:} Jako gracz chcę rozwiązywać różnorodne zagadki na kolejnych poziomach, aby doskonalić umiejętność wykorzystania jednej mechaniki w zmiennych kontekstach.

\textit{Kryteria akceptacji:}
\begin{itemize}[noitemsep]
    \item Mając ukończoną grę, kiedy gracz przeszedł całą sekwencję, wtedy zaliczył co najmniej trzy poziomy zawierające niepowtarzające się zagadki logiczne.
    \item Mając każdy z poziomów, kiedy gracz analizuje jego strukturę, wtedy wymagana do ukończenia interakcja z przeciwnikiem wykorzystuje mechanikę parowania (odbicie pocisku lub zmianę momentum).
\end{itemize}

\subsection{Interakcje środowiskowe}

\subsubsection*{WF-13: Elementy aktywowane parowaniem}
\textbf{User story:} Jako gracz chcę aktywować elementy otoczenia za pomocą odbitych pocisków, aby otwierać drogę do dalszej części poziomu.

\textit{Kryteria akceptacji:}
\begin{itemize}[noitemsep]
    \item Mając element aktywowany na poziomie (np. przełącznik), kiedy zostaje trafiony odbitym pociskiem, wtedy zmienia stan na aktywny.
    \item Mając aktywowany przełącznik powiązany z drzwiami, kiedy zmienia stan na aktywny, wtedy odpowiadające mu drzwi otwierają się.
    \item Mając aktywowany element, kiedy poziom zostaje zrestartowany, wtedy element wraca do stanu początkowego.
\end{itemize}

\subsubsection*{WF-14: Kolizje z otoczeniem}
\textbf{User story:} Jako gracz chcę, aby fizyka kolizji z otoczeniem była przewidywalna, aby planować sekwencje ruchu i odbić.

\textit{Kryteria akceptacji:}
\begin{itemize}[noitemsep]
    \item Mając postać poruszającą się w kierunku obiektu statycznego, kiedy następuje kolizja, wtedy ruch w danym kierunku zostaje zatrzymany bez efektu „wchodzenia w ścianę”.
    \item Mając postać z aktywnym pędem po parowaniu, kiedy następuje kolizja z powierzchnią ustawioną pod kątem, wtedy postać odbija się zgodnie z kątem padania.
    \item Mając postać znajdującą się w powietrzu, kiedy dotyka powierzchni dolnej platformy, wtedy zatrzymuje ruch pionowy i może wykonać kolejny skok.
\end{itemize}

\subsection{Interfejs użytkownika}

\subsubsection*{WF-15: Menu główne}
\textbf{User story:} Jako gracz chcę mieć dostęp do menu głównego, aby rozpocząć grę, kontynuować ostatnio osiągnięty poziom lub zakończyć aplikację.

\textit{Kryteria akceptacji:}
\begin{itemize}[noitemsep]
    \item Mając uruchomioną grę, kiedy gracz uruchamia aplikację, wtedy wyświetlane jest menu główne zawierające co najmniej opcje: \textit{Nowa gra}, \textit{Kontynuuj}, \textit{Wyjście}.
    \item Mając wyświetlone menu główne, kiedy gracz wybiera \textit{Nowa gra}, wtedy uruchamia się pierwszy poziom sekwencji.
    \item Mając wyświetlone menu główne, kiedy gracz wybiera \textit{Kontynuuj}, wtedy uruchamia się ostatnio odblokowany poziom.
\end{itemize}

\subsubsection*{WF-16: Zapis postępu}
\textbf{User story:} Jako gracz chcę, aby gra zapamiętała mój postęp, abym mógł kontynuować rozgrywkę po wyłączeniu aplikacji.

\textit{Kryteria akceptacji:}
\begin{itemize}[noitemsep]
    \item Mając ukończony poziom, kiedy ekran przejścia się zakończył, wtedy informacja o odblokowanym kolejnym poziomie jest zapisywana lokalnie.
    \item Mając wyłączoną i ponownie uruchomioną grę, kiedy gracz wybiera \textit{Kontynuuj}, wtedy uruchamiany jest ostatnio odblokowany poziom.
\end{itemize}

\subsubsection*{WF-17: Ekran końcowy}
\textbf{User story:} Jako gracz chcę otrzymać czytelne podsumowanie po ukończeniu ostatniego poziomu, aby wiedzieć, że gra dobiegła końca.

\textit{Kryteria akceptacji:}
\begin{itemize}[noitemsep]
    \item Mając ukończony ostatni poziom sekwencji, kiedy poziom zostaje zaliczony, wtedy wyświetlany jest ekran zakończenia gry.
    \item Mając wyświetlony ekran końcowy, kiedy gracz potwierdza zakończenie, wtedy następuje powrót do menu głównego.
\end{itemize}

\subsubsection*{WF-18: Pauza}
\textbf{User story:} Jako gracz chcę móc wstrzymać rozgrywkę w dowolnym momencie, aby przerwać sesję bez utraty stanu poziomu.

\textit{Kryteria akceptacji:}
\begin{itemize}[noitemsep]
    \item Mając aktywny poziom, kiedy gracz wciska klawisz pauzy, wtedy symulacja gry zostaje zatrzymana i wyświetlone zostaje menu pauzy.
    \item Mając wyświetlone menu pauzy, kiedy gracz wybiera \textit{Wznów}, wtedy symulacja jest kontynuowana od stanu sprzed pauzy.
    \item Mając wyświetlone menu pauzy, kiedy gracz wybiera \textit{Wyjdź do menu}, wtedy bieżący poziom zostaje porzucony, a gra wraca do menu głównego.
\end{itemize}

\section{Wymagania niefunkcjonalne}

Wymagania niefunkcjonalne określają cechy jakościowe gry, których spełnienie warunkuje grywalność i komfort odbioru. Wszystkie kryteria akceptacji wyrażono w formie mierzalnej i testowalnej. Szczególny nacisk położono na wydajność oraz determinizm symulacji, ponieważ mechanika parowania opiera się na precyzyjnych oknach czasowych liczonych w pojedynczych klatkach (\textit{frame data}); jakiekolwiek odchylenia w tym zakresie podważyłyby fundament rozgrywki.

\subsection{Wydajność}

\subsubsection*{WNF-01: Stabilność liczby klatek}
Gra musi utrzymywać minimalnie 60 klatek na sekundę na sprzęcie referencyjnym (procesor klasy Intel Core i5 8. generacji lub odpowiednik, 8~GB RAM, zintegrowana karta graficzna, dysk SSD) podczas rozgrywki na dowolnym poziomie zawartym w wersji finalnej.

\textit{Sposób weryfikacji:} pomiar FPS narzędziem wbudowanym w silnik gry przez co najmniej 60 sekund rozgrywki na każdym poziomie; dopuszczalne są spadki krótsze niż 100~ms.

\subsubsection*{WNF-02: Opóźnienie wejścia}
Opóźnienie pomiędzy fizycznym naciśnięciem klawisza a reakcją gry (zmiana stanu postaci) nie może przekraczać 50~ms w 95\% pomiarów.

\textit{Sposób weryfikacji:} pomiar narzędziem rejestrującym wejście oraz analiza nagrania ekranu z licznikiem klatek dla 100 prób parowania.

\subsubsection*{WNF-03: Czas restartu poziomu}
Czas pomiędzy śmiercią postaci a ponownym uzyskaniem kontroli przez gracza nie może przekraczać 1 sekundy.

\textit{Sposób weryfikacji:} pomiar dla 30 cykli śmierć--restart na każdym poziomie z agregacją wartości średniej i maksymalnej.

\subsubsection*{WNF-04: Czas ładowania poziomu}
Czas ładowania pojedynczego poziomu nie może przekraczać 3 sekund na sprzęcie referencyjnym.

\textit{Sposób weryfikacji:} pomiar czasu od wyboru poziomu w menu do uzyskania sterowania postaci dla każdego poziomu osobno.

\subsubsection*{WNF-05: Czas uruchomienia gry}
Czas uruchomienia gry od momentu kliknięcia pliku wykonywalnego do wyświetlenia menu głównego nie może przekraczać 30 sekund na sprzęcie referencyjnym.

\textit{Sposób weryfikacji:} pomiar stoperem dla pięciu uruchomień; brany pod uwagę jest wynik średni.

\subsection{Niezawodność}

\subsubsection*{WNF-06: Determinizm symulacji}
Identyczna sekwencja wejść gracza na tym samym poziomie musi prowadzić do identycznego rezultatu fizycznego (pozycja postaci, stan przeciwników).

\textit{Sposób weryfikacji:} odtworzenie zarejestrowanej sekwencji wejść dla 10 wybranych przejść poziomów; porównanie końcowego stanu sceny.

\subsubsection*{WNF-07: Stabilność}
Gra musi działać bez awarii (\textit{crash}) przez co najmniej 4 godziny ciągłej rozgrywki obejmującej powtarzane przejścia poziomów oraz restarty.

\textit{Sposób weryfikacji:} sesja testowa zorganizowana w fazie testów końcowych z udziałem co najmniej dwóch testerów.

\subsubsection*{WNF-08: Zachowanie postępu}
Postęp gracza zapisany lokalnie nie może zostać utracony w wyniku zamknięcia gry przez menu, restartu komputera lub wymuszonego zamknięcia aplikacji po ukończeniu poziomu.

\textit{Sposób weryfikacji:} test scenariuszowy --- ukończenie poziomu, zamknięcie gry trzema sposobami (przez menu, przez \textit{Alt+F4}, przez wymuszone zakończenie procesu), weryfikacja stanu zapisu po ponownym uruchomieniu.

\subsection{Użyteczność}

\subsubsection*{WNF-09: Czytelność informacji zwrotnej}
Każde udane oraz nieudane parowanie musi być sygnalizowane efektem wizualnym oraz dźwiękowym rozróżnialnym w czasie krótszym niż 100~ms od zdarzenia.

\textit{Sposób weryfikacji:} test obserwacyjny z udziałem co najmniej trzech osób spoza zespołu projektowego; każda z nich powinna prawidłowo rozróżnić wszystkie typy informacji zwrotnej.

\subsubsection*{WNF-10: Krzywa uczenia się}
Osoba niezaznajomiona z grą musi być w stanie ukończyć pierwszy poziom w czasie nieprzekraczającym 10 minut bez konieczności konsultowania zewnętrznych instrukcji.

\textit{Sposób weryfikacji:} playtest z udziałem co najmniej trzech osób spoza zespołu projektowego; pomiar czasu ukończenia oraz zliczenie zadanych pytań.

\subsubsection*{WNF-11: Konfigurowalność sterowania}
Domyślny układ sterowania musi być zgodny z konwencjami gatunku (klawisze WASD lub strzałki dla ruchu, Spacja dla skoku, jeden dedykowany klawisz dla parowania). W wersji finalnej dopuszczalne jest sterowanie wyłącznie z klawiatury.

\textit{Sposób weryfikacji:} przegląd manualny; weryfikacja zgodności z opisem w dokumentacji.

\subsubsection*{WNF-12: Czytelność elementów rozgrywki}
Wszystkie obiekty interaktywne (przeciwnicy, pociski, przełączniki, punkt końcowy poziomu) muszą być wizualnie odróżnialne od obiektów statycznych otoczenia oraz wzajemnie między sobą.

\textit{Sposób weryfikacji:} ocena ekspercka przez Głównego Użytkownika oraz potwierdzenie w playtestach (brak pytań typu ,,czego mam dotknąć?'' podczas testów).

\subsection{Kompatybilność i wdrożenie}

\subsubsection*{WNF-13: Platforma docelowa}
Gra musi uruchamiać się i działać poprawnie na systemie operacyjnym Windows 10 oraz Windows 11 w wersji 64-bitowej.

\textit{Sposób weryfikacji:} test uruchomienia oraz przejście co najmniej jednego poziomu na każdej z platform.

\subsubsection*{WNF-14: Rozdzielczość ekranu}
Gra musi obsługiwać natywnie rozdzielczość 1920$\times$1080 oraz skalować interfejs do rozdzielczości od 1280$\times$720 do 2560$\times$1440 bez utraty czytelności i bez przycięć elementów HUD.

\textit{Sposób weryfikacji:} test wizualny w trzech rozdzielczościach skrajnych i jednej pośredniej.

\subsubsection*{WNF-15: Rozmiar dystrybucji}
Rozmiar wersji dystrybucyjnej gry (plik wykonywalny wraz z zasobami) nie może przekraczać 15~GB.

\textit{Sposób weryfikacji:} pomiar rozmiaru folderu wersji wydanej.

\subsection{Utrzymywalność i rozszerzalność}

\subsubsection*{WNF-16: Struktura danych poziomów}
Poziomy gry muszą być przechowywane w plikach o formacie czytelnym dla człowieka (\textit{plain text} lub format silnika), umożliwiającym dodanie kolejnego poziomu bez modyfikacji kodu źródłowego logiki gry.

\textit{Sposób weryfikacji:} dodanie poziomu testowego do gry z poziomu pliku danych i weryfikacja jego poprawnego załadowania.

\subsubsection*{WNF-17: Parametryzacja przeciwników}
Kluczowe parametry przeciwników (interwał ataku, prędkość pocisku, długość fazy zapowiedzi) muszą być definiowane w plikach konfiguracyjnych poza kodem źródłowym.

\textit{Sposób weryfikacji:} zmiana parametru w pliku konfiguracyjnym i weryfikacja wpływu na zachowanie przeciwnika bez konieczności rekompilacji.

\subsubsection*{WNF-18: Repozytorium kodu}
Cały kod źródłowy gry musi być przechowywany w repozytorium z historią zmian umożliwiającą identyfikację autora oraz daty każdej zmiany.

\textit{Sposób weryfikacji:} weryfikacja stanu repozytorium na koniec każdego sprintu w ramach \textit{code review}.


\end{document}